\documentclass[11pt]{article}
\usepackage[a4paper,margin=25mm]{geometry}
\usepackage{amsmath,amssymb,amsthm,booktabs,array}
\usepackage[T1]{fontenc}
\usepackage{lmodern}
\usepackage[hidelinks]{hyperref}
\newtheorem{theorem}{Theorem}
\newtheorem{corollary}[theorem]{Corollary}
\newcommand{\ind}{\mathbf{1}}
\title{Extensions of Shrock's and Yurishchev's formulas\protect\\to Hardin's defective-colouring recurrences}
\author{Haobo Ma\thanks{ChronoAI Pte Ltd, Singapore.}\and
Wenlin Zhang\thanks{National University of Singapore; ChronoAI Pte Ltd, Singapore.}}
\date{19 November 2026}
\begin{document}
\maketitle
\begin{abstract}
Extending Shrock's open-ladder formula (2000,
Eqs.~(5.1)--(5.6) \cite{shrock}) and Yurishchev's double-chain block
(2001, Eq.~(9) \cite{yurishchev}), an exact defect-counting transfer
proves R. H. Hardin's empirical recurrence in OEIS A229634, Formula,
revision 9 (27 April 2021; first posted 27 September 2013),
for every $n\geq21$.
Its order eighteen and its starting index are sharp.
The same construction proves every displayed row and column recurrence
of the companion table A229637, with minimal eventual orders and sharp
starting indices. A general defect-budget argument explains the cubed
denominators. The two-column formulas of Colin Barker and the
three- and four-column generating functions of R. J. Mathar and Barker
follow as well.
\end{abstract}

\section{The count and the recurrence}
An $n\times k$ array has entries in $\{0,1,2\}$.
Its bonds join horizontally consecutive cells in a row, and diagonally
consecutive cells in consecutive rows, in both diagonal directions.
There are no vertical bonds. A mistake is a bond whose endpoints have
the same colour. Colours are introduced in row-major order: the first
colour is $0$, and the first occurrence of a new colour uses the smallest
unused label. Write $a_k(n)$ for the number of such canonical arrays
with exactly two mistakes, and put $a(n)=a_5(n)$.
This is Hardin's A229634, the fifth column of A229637
\cite{hardin5,hardintable}. In particular, the source calls the proposed
formula ``Empirical'', rather than a proved recurrence.

Define $c_0,\ldots,c_{17}$, in increasing order of the lagged index, by
\begin{equation}\label{eq:coefficients}
\begin{split}
(c_0,\ldots,c_{17})={}&(8,-228,2598,-14887,42804,-40719,\\
&-84933,243513,-118011,-230994,305043,-55215,\\
&-116010,94839,-33687,6569,-723,42).
\end{split}
\end{equation}
\begin{theorem}\label{thm:main}
For every integer $n\geq21$,
\begin{equation}\label{eq:recurrence}
a(n)=\sum_{j=0}^{17}c_j a(n-18+j).
\end{equation}
This is exactly the recurrence in the Formula field of A229634.
\end{theorem}

We prove the theorem by a finite vector identity and a propagation
argument. The identity is stronger than agreement of a finite list of
terms: it holds for the whole distribution of the last row and unused
defect budget.

\section{An exact transfer proof}
Let $R_k=\{0,1,2\}^k$. For rows $p,r\in R_k$, define
\[
h(r)=\sum_{i=1}^{k-1}\ind_{r_i=r_{i+1}},\qquad
b(p,r)=\sum_{i=1}^{k-1}
  (\ind_{p_i=r_{i+1}}+\ind_{p_{i+1}=r_i}).
\]
Thus $h$ counts horizontal mistakes and $b$ counts mistakes between
two adjacent rows. Use states $(r,d)\in R_k\times\{0,1,2\}$, where $d$
is the remaining number of mistakes. The integer matrix $M_k$ is
defined by
\[
(M_k)_{(p,d),(r,e)}=
\begin{cases}
1,&e=d-h(r)-b(p,r)\ \text{and}\ e\geq0,\\
0,&\text{otherwise}.
\end{cases}
\]
Let $w_0$ be the row vector with
$w_0(r,d)=\ind_{d=2-h(r)}$, and set $w_t=w_0M_k^t$.
Finally, let $v(r,d)=\ind_{d=0}$ be a column vector.
The scalar $L_k(n)=w_{n-1}v$ counts all labelled arrays with exactly
two mistakes. Indeed, initialization charges the first row's horizontal
bonds, each transition charges precisely the newly added horizontal and
diagonal bonds, and $v$ accepts precisely the exhausted budgets.
This gives a bijection between counted arrays and accepted paths;
budgets never increase, so no rejected path can later become admissible.

For $k=5$, every counted array is nonconstant. A constant array would
have $4n+8(n-1)=12n-8$ mistakes, which is never two.
A nonconstant array uses at least two colours, so its stabilizer under
the six colour permutations is trivial. Each orbit has exactly one
row-major canonical representative, obtained by relabelling colours
in their order of first occurrence. Consequently
\begin{equation}\label{eq:orbit}
L_5(n)=6a(n)\qquad(n\geq1).
\end{equation}

For the explicit matrix and vector above, the finite identity is
\begin{equation}\label{eq:vector}
w_{20}=\sum_{j=0}^{17}c_j w_{2+j}\qquad(k=5).
\end{equation}
Here is a complete specification of its exact-arithmetic calculation.
Order the rows lexicographically, order the remaining budgets as $0,1,2$,
initialize $w_0$ as above, and repeat
\[
w_{t+1}(r,e)=
\sum_{\substack{p\in R_5,\ d\in\{0,1,2\}\\
 e=d-h(r)-b(p,r)\geq0}} w_t(p,d)
\quad(0\leq t<20).
\]
Subtract the right side of \eqref{eq:vector} coordinate by coordinate;
every resulting integer is zero. The accompanying \texttt{verify.py}
implements precisely these sums, with integer arithmetic and no
preassigned transition table. Its output gives the coefficient list,
the initial scalar values, and the finite identities used here.

Multiplication of \eqref{eq:vector} by $M_5^m v$, for any $m\geq0$,
gives
\[
L_5(21+m)=\sum_{j=0}^{17}c_jL_5(3+j+m).
\]
Dividing by six using \eqref{eq:orbit} proves Theorem~\ref{thm:main}.
\qed

\section{The generating function and optimal order}
Put $a(0)=0$ and $A(z)=\sum_{n\geq0}a(n)z^n$. Define
\[
Q(z)=(1-14z+45z^2-15z^3-36z^4+19z^5-2z^6)^3.
\]
Expansion gives $Q(z)=1-\sum_{j=0}^{17}c_jz^{18-j}$.
The initial transfer values and Theorem~\ref{thm:main} give
\begin{equation}\label{eq:gf}
A(z)=\frac{P(z)}{Q(z)},\qquad
P(z)=\sum_{i=1}^{20}p_i z^i,
\end{equation}
where the coefficients are
\begin{center}
\begin{tabular}{rr@{\qquad}rr@{\qquad}rr@{\qquad}rr}
\toprule
$i$&$p_i$&$i$&$p_i$&$i$&$p_i$&$i$&$p_i$\\
\midrule
1&12&6&2954875&11&2209941&16&$-1304212$\\
2&$-16$&7&$-5909268$&12&11838704&17&320998\\
3&$-5789$&8&2800838&13&$-8530008$&18&$-42652$\\
4&98314&9&10098304&14&$-59389$&19&2920\\
5&$-740730$&10&$-16241498$&15&2512320&20&$-80$\\
\bottomrule
\end{tabular}
\end{center}
For completeness, the computation of each coefficient is simply
$p_i=\sum_{\ell=0}^{\min(i,18)}[z^\ell]Q(z)\,a(i-\ell)$;
all higher coefficients vanish by the recurrence.

\begin{corollary}\label{cor:minimal}
The least order of an eventual rational constant-coefficient recurrence
for $a(n)$ is eighteen. Equation~\eqref{eq:recurrence} fails at $n=20$.
\end{corollary}
\begin{proof}
The polynomials $P$ and $Q$ are coprime. One finite certificate is their
Euclidean algorithm modulo $101$: the successive nonzero remainder
degrees are $17,16,\ldots,0$, with last constant $63$.
Both leading coefficients remain nonzero modulo $101$.
Hence their resultant is nonzero modulo $101$, and therefore nonzero
over $\mathbb Q$. These polynomial computations are also explicit in
\texttt{verify.py}.
If an eventual recurrence has denominator polynomial $D(z)$, then
$D(z)A(z)$ is a polynomial. Equation~\eqref{eq:gf} and coprimality
imply $Q\mid D$, so its order is at least $\deg Q=18$.
Theorem~\ref{thm:main} attains this order.
Finally, $[z^{20}]Q(z)A(z)=p_{20}=-80\ne0$, which is exactly the
residual $a(20)-\sum_{j=0}^{17}c_j a(2+j)$.
\end{proof}

\section{Why the denominators are cubes}
Transfer matrices for strip colourings are classical; see Salas and
Sokal \cite[Section 3]{salassokal}. Here the defect budget gives an
additional denominator bound.
Let a slice have $s$ sites and $q\geq1$ colours. Let $h(r)$ and
$b(p,r)$ be fixed nonnegative integer costs within a slice and between
consecutive slices. For $d\geq0$, let $L_d(\ell)$ count labelled
strings of $\ell\geq1$ slices with total cost exactly $d$.
Put $T_0(p,r)=\ind_{h(r)+b(p,r)=0}$ and $D(z)=\det(I-zT_0)$.
\begin{theorem}\label{thm:budget}
The reduced denominator of $\sum_{\ell\geq1}L_d(\ell)z^\ell$ divides
$D(z)^{d+1}$. If a group acting on slices preserves $h$ and $b$, the
same assertion holds with the orbit transition matrix replacing $T_0$.
\end{theorem}
\begin{proof}
Use states $(r,e)$ with remaining budget $0\leq e\leq d$, initial
weight $\ind_{e=d-h(r)}$, and terminal weight $\ind_{e=0}$.
The transition decreases $e$ by $h(r)+b(p,r)$. Ordering states by
budget makes the matrix block triangular, with $d+1$ identical
diagonal blocks $T_0$. Thus $\det(I-zM)=D(z)^{d+1}$.
The path interpretation gives the generating function
$z w_0(I-zM)^{-1}v$; the adjugate formula proves the assertion.
For an orbit $C$, the number of $r\in C$ with each prescribed cost
is independent of the representative $p$, because the group preserves
costs. Use these numbers as transition weights and sum initial weights
over orbits. Induction preserves total weight in every orbit and
budget. The diagonal blocks are again identical, proving the assertion.
\end{proof}

For the present arrays, use the group generated by colour permutations
and slice reversal. Represent each orbit by the lexicographically
smaller of the two row-major canonical labellings of the slice and its
reversal. Reversal groups transfer states; it does not identify arrays
in the final count. For fixed width $k$, use the costs in Section 2.
For fixed height $n$, use column slices and
\[
h(r)=0,\qquad b(p,r)=\sum_{i=1}^{n}\ind_{p_i=r_i}
+\sum_{i=1}^{n-1}(\ind_{p_i=r_{i+1}}+\ind_{p_{i+1}=r_i}).
\]
The first sum counts horizontal bonds. Denote the orbit zero-cost
matrix by $B$, and put $D=\det(I-zB)$. For fixed width, columns
indexed by slices with $h(r)>0$ are zero in $B$; deleting their rows
and columns leaves this determinant unchanged. These rules specify
every entry without a preassigned transition table.

\section{All displayed companion recurrences}
Let $T(n,k)=a_k(n)$. Use $x=n$ in a fixed column and $x=k$ in a
fixed row, and set each sequence to zero at index zero.
For $Q(z)=D(z)^3=\sum_{j=0}^{r}q_jz^j$, with $q_0=1$, its
associated recurrence means
\begin{equation}\label{eq:allrec}
f(x)=-\sum_{j=1}^{r}q_j f(x-j).
\end{equation}
\begin{corollary}\label{cor:all}
Every displayed empirical recurrence in A229637, Formula, revision 8
(27 April 2021), holds. In the following table $r$ is the least
eventual rational constant-coefficient order, and $N$ is the sharp
starting index of \eqref{eq:allrec}. The finite vector identity below
has index $N_v$. The last column is the residual at $x=N-1$.
\end{corollary}
\begin{center}\small
\begin{tabular}{lrrrr}
\toprule
Sequence&$r$&$N$&$N_v$&Residual\\
\midrule
$T(x,2)$&3&6&6&$-2$\\
$T(x,3)$&6&8&8&$-88$\\
$T(x,4)$&6&10&11&25\\
$T(x,5)$&18&21&21&$-80$\\
$T(x,6)$&27&31&32&128\\
$T(x,7)$&57&61&62&$-20736$\\
\midrule
$T(2,x)$&6&7&7&$-20$\\
$T(3,x)$&9&13&13&8\\
$T(4,x)$&18&22&22&$-1280$\\
$T(5,x)$&30&34&34&8211456\\
$T(6,x)$&69&73&73&88260480000\\
\bottomrule
\end{tabular}
\end{center}
\begin{proof}
Use the orbit transfer above with budget two, initializing
$w_0(C,e)=|C|\ind_{e=2-h(C)}$. Its terminal weight is the labelled
array count $L$. A constant $n\times k$ array has $(k-1)(3n-2)$
mistakes. Thus the canonical count is $L/6$, except at $(1,3)$,
where it is $(L+3)/6=1$. All other counted arrays have colour orbits
of size six, as in Section 2.

For each table row, the exact finite vector identity is
\[
\sum_{j=0}^{r}q_j w_{N_v-1-j}=0.
\]
The accompanying \texttt{verify\_siblings.py} constructs all orbit
transitions from the three-colour slices, computes $D=\det(I-zB)$
by integer characteristic-polynomial arithmetic, and verifies this
identity in every coordinate. Subsequent transfer powers give
\eqref{eq:allrec} for $x\geq N_v$.
The exceptional constant array $(1,3)$ occurs in none of these identities.
Where $N_v=N+1$, the scalar residual at $N$ is also zero, establishing
the stated range. These are finite vector identities rather than an
extrapolation from an initial list of terms.

The file \texttt{sibling-output.json} supplies $D,Q,P$ in ascending
coefficient order, where
$P_x=\sum_{j=0}^{\min(x,r)}q_jf(x-j)$ for $x<N$.
It verifies the displayed $P_{N-1}$; subsequent coefficients vanish
by the recurrence. For every row, the Euclidean algorithm modulo
$101$ ends in a nonzero constant, with neither leading coefficient
zero. Thus $\gcd(P,Q)=1$ over $\mathbb Q$.
As in Corollary~\ref{cor:minimal}, the denominator has minimal degree
$r$, and $P_{N-1}\ne0$ makes $N$ sharp.
\end{proof}

For orientation, the first four fixed-column polynomials are
\[
\begin{array}{ll}
k=2:&D=1-3z,\\
k=3:&D=1-5z+2z^2,\\
k=4:&D=1-7z+4z^2,\\
k=5:&D=1-14z+45z^2-15z^3-36z^4+19z^5-2z^6.
\end{array}
\]
For row two, $D=1-2z-2z^2$; cubing it gives exactly the six
coefficients in A229637. For columns four, six and seven and rows
three through six, that entry specifies only orders, not coefficients;
the determinant construction specifies the recurrences explicitly.

\begin{corollary}\label{cor:barker}
The two-column sequence is A229600, and
\[
\sum_{n\geq1}a_2(n)z^n
=\frac{z^2(6-15z+13z^2-2z^3)}{(1-3z)^3},\qquad
a_2(n)=3^{n-5}(32n^2+31n-30)\quad(n\geq3).
\]
These are Barker's conjectures in A229600, Formula, 16 March 2018.
They also follow from Shrock's earlier open-ladder partition function
\cite[Eqs.~(5.1)--(5.6)]{shrock} by setting $q=3$, $v=t-1$ and
taking the coefficient of $t^2$: precisely,
$\sum_{n\geq1}a_2(n)z^n=z^2[t^2]\Gamma(S,3,t-1,z)/6$.
The empirical generating functions credited to Mathar in A229632
(29 September 2013) and Barker in A229633 (20 September 2018)
also hold \cite{widthtwo,widththree,widthfour}.
\end{corollary}
\begin{proof}
Swap the two cells in each even-numbered row. Horizontal and crossed
diagonal bonds become the horizontal and vertical bonds of A229600's
ladder. This colour-equivariant graph isomorphism gives a bijection
on canonical colour orbits, even though individual row-major labels
may change. The numerator in Corollary~\ref{cor:all} is
$6z^2-15z^3+13z^4-2z^5$.
Using $[z^m](1-3z)^{-3}=\binom{m+2}{2}3^m$ yields the stated
formula for $n\geq5$; it also gives $39,202$ at $n=3,4$.
For width three, its numerator is
$zD^3-4z^2(2z-5)(12z^4-45z^3+47z^2-17z+2)$.
For width four, it is
$z(3+59z-976z^2+4803z^3-9539z^4+7326z^5-1825z^6-76z^7+25z^8)$.
These are precisely the two cited generating functions.
\end{proof}

The two-row case has a separate exact antecedent. Specialize
Yurishchev's block \cite[Eqs.~(5), (8), (9)]{yurishchev} to three
colours, give the four intercolumn bonds weight $t$, and set all
other couplings to zero. The normalized equal-colour and distinct-colour
slice sums give
\[
H(t)=\begin{pmatrix}
2+t^4&\sqrt2(1+2t^2)\\
\sqrt2(1+2t^2)&4t+2t^2
\end{pmatrix},\qquad b=\begin{pmatrix}\sqrt3\\\sqrt6\end{pmatrix}.
\]
For example, from an equal-colour slice there is one equal-colour
successor of weight $t^4$ and two of weight one. For a fixed
equal-colour slice, two distinct-colour successors have weight one and
four have weight $t^2$; normalization divides their sum by $\sqrt2$.
From a distinct-colour slice,
four distinct successors have one equal bond and two have two.
The free-boundary labelled polynomial is $b^{\mathsf T}H(t)^{k-1}b$:
the all-ones boundary vector has coordinates $b$ in this normalized
invariant basis. Summing the geometric series and taking its $t^2$
coefficient gives
\[
\sum_{k\geq1}T(2,k)z^k
=\frac16[t^2]z b^{\mathsf T}(I-zH(t))^{-1}b
=\frac{2z^2(1-z)(3+5z-36z^2+10z^3)}{(1-2z-2z^2)^3}.
\]
This is also Barker's empirical generating function in A229638,
Formula, 20 September 2018 \cite{rowtwo}. The free-boundary
contraction and coefficient extraction above make the connection
explicit; the published block, rather than this particular finite
strip formula, is Yurishchev's result.

\begin{corollary}\label{cor:rowone}
One has $T(n,1)=0$ for $n\geq1$ and
\[
T(1,k)=\binom{k-1}{2}2^{k-4}\quad(k\geq4),\qquad
(T(1,1),T(1,2),T(1,3))=(0,0,1).
\]
Consequently $T(1,k)=6T(1,k-1)-12T(1,k-2)+8T(1,k-3)$ for
$k\geq7$, with residual $-4$ at $k=6$ and minimal eventual order three.
For $k\geq4$ this is Sloane's A052482$(k-2)$ \cite{sloane}.
In this notation the reciprocal sums recorded there by Amiram Eldar
(28 September 2025) are
\[
\sum_{k\geq4}\frac1{T(1,k)}=6-8\log2,\qquad
\sum_{k\geq4}\frac{(-1)^k}{T(1,k)}=10-24\log(3/2).
\]
\end{corollary}
\begin{proof}
Width one has no bonds. For one row of width $k\geq4$, choose two
equal bonds among $k-1$, then choose the first colour in three ways
and each unequal bond in two ways. Divide by six. The smaller values
follow directly. The third difference of $T(1,k)/2^k$ is zero for
$k\geq7$; its nonzero quadratic term makes the reduced denominator
$(1-2z)^3$, so the order is minimal. Substitution gives the residual.
For the sums, put $m=k-2$ and subtract the logarithmic series for
$1/m$ and $1/(m+1)$ to obtain
\[
S(x)=\sum_{m\geq2}\frac{x^m}{m(m+1)}
=1-\frac{x}{2}+\left(\frac1x-1\right)\log(1-x),\quad 0<|x|<1.
\]
Evaluate $8S(1/2)$ and $8S(-1/2)$.
\end{proof}

\begin{corollary}\label{cor:zero}
Let $U(n)$ be the number of labelled $n\times5$ arrays in A232917:
no horizontal or diagonal neighbours have values summing to two.
Then
\[
\sum_{n\geq1}U(n)z^n
=\frac{6z(4-7z+z^2)(2-8z-5z^2+3z^3)}
 {1-14z+45z^2-15z^3-36z^4+19z^5-2z^6}.
\]
This proves Barker's empirical generating function (6 October 2018)
and Hardin's order-six recurrence in A232917 \cite{zerowidthfive},
the latter for $n\geq7$.
\end{corollary}
\begin{proof}
Replace each colour $x$ by $2-x$ in alternate columns. Every bond
changes column parity, so a complementary pair becomes an equal pair,
and conversely. This bijects the arrays with the labelled zero-defect
arrays of Section 2. The zero-cost orbit matrix $B$ has six states,
so $z u(I-zB)^{-1}\mathbf1$ has denominator $D$ and numerator of
degree at most six, where $u(C)=|C|$ for proper slices.
The first six zero-defect counts, calculated by the same transitions,
give the displayed numerator. The accompanying checker verifies its
coefficients exactly. Expanding $D$ gives the stated recurrence.
\end{proof}

\begin{center}\small
\begin{tabular}{>{\raggedright\arraybackslash}p{.43\linewidth}>{\raggedright\arraybackslash}p{.49\linewidth}}
\toprule
Source item&Coverage\\
\midrule
A229634, Formula; A229637, column 5&Proved; sharp range and minimal order.\\
A229637, columns 2 and row 2; A229600 and A229638&Reproved from prior
formulas: Shrock (2000), Yurishchev (2001); explicit extractions above.\\
A229637, remaining displayed columns and rows&Proved;
Corollaries~\ref{cor:all} and \ref{cor:rowone}.\\
A229632 and A229633, Formula&Hardin's recurrences and Mathar's and Barker's
empirical generating functions proved.\\
A052482, Name; Barker's generating function and recurrence&Prior formulas cited
(Sloane, 2000; Barker, 2017); one-row interpretation proved.\\
A052482, reciprocal sums&Prior formulas cited (Eldar, 2025), with the deductions above.\\
A232917, empirical recurrence and Barker generating function&Proved;
Corollary~\ref{cor:zero}; its denominator is the five-column zero-defect determinant.\\
A052482, folded-cube interpretations; A229606, other widths&Not covered.\\
\bottomrule
\end{tabular}
\end{center}


The pinned Lean declaration
\href{https://github.com/the-omega-institute/trureturing/blob/9849c0d81bfe0a953a10ad37400f5b1b61f028b5/D5/S3/StatisticalMechanics/Potts/CrossedStripDefectsWide/HardinA229634Strip.lean}{\texttt{HardinA229634Strip.result}}
states the recurrence for the literal canonical count as
$a(n+21)=\sum_{j=0}^{17}c_j a(n+3+j)$ for every natural $n$.
The pinned declarations
\href{https://github.com/the-omega-institute/trureturing/blob/9849c0d81bfe0a953a10ad37400f5b1b61f028b5/D5/S3/StatisticalMechanics/Potts/CrossedStripDefects/HardinA229632Strip.lean}{\texttt{HardinA229632Strip.result}}
and
\href{https://github.com/the-omega-institute/trureturing/blob/9849c0d81bfe0a953a10ad37400f5b1b61f028b5/D5/S3/StatisticalMechanics/Potts/CrossedStripDefects/HardinA229633Strip.lean}{\texttt{HardinA229633Strip.result}}
state the literal width-three and width-four recurrences for $n\geq8$
and $n\geq10$, respectively. The generating functions, optimal orders,
boundary residuals, general budget theorem and other companion results
have the written proofs and exact finite checks specified here;
they are separate from these declarations' scope.
The pinned declaration
\href{https://github.com/the-omega-institute/trureturing/blob/9849c0d81bfe0a953a10ad37400f5b1b61f028b5/D5/S3/StatisticalMechanics/Potts/StripDefects/HardinA229638Strip.lean}{\texttt{HardinA229638Strip.result}}
states precisely the two-row recurrence for $k\geq7$; its generating
function and optimality assertions have the separate proofs above.

\begin{thebibliography}{99}
\bibitem{hardin5}
R. H. Hardin, \emph{Number of defective 3-colorings of an $n\times5$
0..2 array connected horizontally, diagonally and antidiagonally with
exactly two mistakes, and colors introduced in row-major 0..2 order},
OEIS A229634, 27 September 2013, revision 9, 27 April 2021.
\url{https://oeis.org/A229634}.

\bibitem{hardintable}
R. H. Hardin, \emph{Defective 3-colorings of an $n\times k$ array
with exactly two mistakes}, OEIS A229637, 27 September 2013,
revision 8, 27 April 2021. \url{https://oeis.org/A229637}.

\bibitem{sloane}
N. J. A. Sloane, \emph{$a(n)=2^{n-2}\binom{n+1}{2}$},
OEIS A052482, 16 March 2000; Name and Formula, revision 20,
28 September 2025, including contributions of C. Barker
(22 September 2017) and A. Eldar (28 September 2025).
\url{https://oeis.org/A052482}.

\bibitem{widthtwo}
R. H. Hardin, OEIS A229600, 26 September 2013;
C. Barker, conjectures in Formula, 16 March 2018.
\url{https://oeis.org/A229600}.

\bibitem{widththree}
R. H. Hardin, OEIS A229632, 27 September 2013;
R. J. Mathar, empirical generating function in Formula, 29 September 2013.
\url{https://oeis.org/A229632}.

\bibitem{widthfour}
R. H. Hardin, OEIS A229633, 27 September 2013;
C. Barker, empirical generating function in Formula, 20 September 2018.
\url{https://oeis.org/A229633}.

\bibitem{salassokal}
J. Salas and A. D. Sokal, \emph{Transfer matrices and partition-function zeros
for antiferromagnetic Potts models. I. General theory and square-lattice
chromatic polynomial}, J. Statist. Phys. \textbf{104} (2001), 609--699,
\href{https://doi.org/10.1023/A:1010376605067}{doi:10.1023/A:1010376605067};
\href{https://arxiv.org/abs/cond-mat/0004330v2}{arXiv:cond-mat/0004330v2}, Section 3.
\bibitem{shrock}
R. Shrock, \emph{Exact Potts model partition functions on ladder graphs},
Physica A \textbf{283} (2000), 388--446,
\href{https://doi.org/10.1016/S0378-4371(00)00109-6}{doi:10.1016/S0378-4371(00)00109-6};
\href{https://arxiv.org/abs/cond-mat/0001389v1}{arXiv:cond-mat/0001389v1},
Section V, Eqs.~(5.1)--(5.6), printed pp.~22--23.

\bibitem{yurishchev}
M. A. Yurishchev, \emph{Double Potts chain and exact results for some
two-dimensional spin models}, J. Exp. Theor. Phys. \textbf{93} (2001),
1113--1118, \href{https://doi.org/10.1134/1.1427183}{doi:10.1134/1.1427183};
\href{https://arxiv.org/abs/cond-mat/0111401v1}{arXiv:cond-mat/0111401v1},
Section 2, Eqs.~(5), (8), (9), printed pp.~4--5.

\bibitem{rowtwo}
R. H. Hardin, OEIS A229638, 27 September 2013;
C. Barker, empirical generating function in Formula, 20 September 2018.
\url{https://oeis.org/A229638}.

\bibitem{zerowidthfive}
R. H. Hardin, OEIS A232917, 2 December 2013;
C. Barker, empirical generating function in Formula, 6 October 2018.
\url{https://oeis.org/A232917}.
\end{thebibliography}
\noindent\small\copyright\ 2026 Haobo Ma and Wenlin Zhang.
Licensed under Creative Commons Attribution 4.0 International
(CC BY 4.0).
\end{document}
